\section{Principal components}
\label{sec:pca}

\subsection{Setup: linear combinations of the observed variables}
\label{subsec:pca_setup}

Let the random vector $\vb{X}^\tp=(X_1,X_2,\dots,X_p)$ have covariance matrix
$\Si$ with eigenvalues $\la_1\ge\la_2\ge\dots\ge\la_p\ge 0$. Consider the linear
combinations

\begin{equation}
  \label{eq:pca_linear_combinations}
  \small
  \begin{aligned}
    Y_1 &= \vb{a}_1^\tp\vb{X} = a_{11}X_1 + a_{12}X_2 + \dots + a_{1p}X_p,\\
    Y_2 &= \vb{a}_2^\tp\vb{X} = a_{21}X_1 + a_{22}X_2 + \dots + a_{2p}X_p,\\
        &\;\;\vdots\\
    Y_p &= \vb{a}_p^\tp\vb{X} = a_{p1}X_1 + a_{p2}X_2 + \dots + a_{pp}X_p.
  \end{aligned}
\end{equation}

Using the bilinearity of the covariance operator we obtain
\begin{align}
  \Var(Y_i)     &= \vb{a}_i^\tp\Si\,\vb{a}_i,
  \label{eq:var_Yi}\\
  \Cov(Y_i,Y_k) &= \vb{a}_i^\tp\Si\,\vb{a}_k,
  \label{eq:cov_YiYk}
\end{align}
for $i,k=1,2,\dots,p$.

\begin{definition}[Principal components]
  The \textbf{principal components} are those uncorrelated linear combinations
  $Y_1,Y_2,\dots,Y_p$ of \eqref{eq:pca_linear_combinations} whose variances
  \eqref{eq:var_Yi} are as large as possible.
\end{definition}

The first linear combination is the one of maximum variance, i.e.\ the one
maximizing $\Var(Y_1)=\vb{a}_1^\tp\Si\vb{a}_1$. Since this variance can be made
arbitrarily large by rescaling $\vb{a}_1$, the indeterminacy is eliminated by
restricting attention to coefficient vectors of unit length:
\begin{equation}
  \label{eq:first_pc_problem}
  \begin{gathered}
    \text{first principal component}=\vb{a}_1^\tp\vb{X}\\
    \text{that maximizes }\;\Var(\vb{a}_1^\tp\vb{X})
    \;\text{ s.t. }\;\vb{a}_1^\tp\vb{a}_1=1.
  \end{gathered}
\end{equation}

In general, the $i$-th principal component is the combination
$\vb{a}_i^\tp\vb{X}$ that maximizes $\Var(\vb{a}_i^\tp\vb{X})$ subject to
$\vb{a}_i^\tp\vb{a}_i=1$ and
$\Cov(\vb{a}_i^\tp\vb{X},\vb{a}_k^\tp\vb{X})=0$ for all $k<i$.

\subsection{The maximization lemma and the solution}
\label{subsec:pca_maximization}

The constrained problem \eqref{eq:first_pc_problem} is solved by the
maximization of a Rayleigh quotient.

\begin{lemma}[Maximization of a quadratic form]
  Let $\mat{B}=\Si$ be positive definite with eigenvalue--eigenvector pairs
  $(\la_1,\vb{e}_1),\dots,(\la_p,\vb{e}_p)$ and $\la_1\ge\dots\ge\la_p\ge0$.
  Then
  \begin{equation}
    \label{eq:rayleigh}
    \max_{\vb{a}\neq\vb{0}}
      \frac{\vb{a}^\tp\Si\vb{a}}{\vb{a}^\tp\vb{a}}=\la_1,
  \end{equation}
  attained at $\vb{a}=\vb{e}_1$, and, restricting to the orthogonal complement
  of the leading eigenvectors,
  \begin{equation}
    \label{eq:rayleigh_k}
    \max_{\substack{\vb{a}\,\perp\,\vb{e}_1,\dots,\vb{e}_k}}
      \frac{\vb{a}^\tp\Si\vb{a}}{\vb{a}^\tp\vb{a}}=\la_{k+1},
  \end{equation}
  attained at $\vb{a}=\vb{e}_{k+1}$, for $k=1,2,\dots,p-1$.
\end{lemma}

\begin{proof}
  Since the eigenvectors are normalized, $\vb{e}_1^\tp\vb{e}_1=1$, and
  evaluating the quotient \eqref{eq:rayleigh} at $\vb{a}=\vb{e}_1$ gives
  \[
    \max_{\vb{a}\neq\vb{0}}
      \frac{\vb{a}^\tp\Si\vb{a}}{\vb{a}^\tp\vb{a}}
    =\la_1
    =\frac{\vb{e}_1^\tp\Si\vb{e}_1}{\vb{e}_1^\tp\vb{e}_1}
    =\vb{e}_1^\tp\Si\vb{e}_1
    =\Var(Y_1).
  \]
  Similarly, for the choice $\vb{a}=\vb{e}_{k+1}$, which satisfies
  $\vb{e}_{k+1}^\tp\vb{e}_i=0$ for $i=1,2,\dots,k$ and $k=1,2,\dots,p-1$,
  \[
    \frac{\vb{e}_{k+1}^\tp\Si\vb{e}_{k+1}}
         {\vb{e}_{k+1}^\tp\vb{e}_{k+1}}
    =\vb{e}_{k+1}^\tp\Si\vb{e}_{k+1}=\Var(Y_{k+1}),
  \]
  and since $\Si\vb{e}_{k+1}=\la_{k+1}\vb{e}_{k+1}$,
  \[
    \vb{e}_{k+1}^\tp(\Si\vb{e}_{k+1})
    =\la_{k+1}\,\vb{e}_{k+1}^\tp\vb{e}_{k+1}=\la_{k+1},
  \]
  so that $\Var(Y_{k+1})=\la_{k+1}$.
\end{proof}

\begin{theorem}[Principal components of $\Si$]
  \label{thm:pca_main}
  Let $\Si$ be the covariance matrix associated with the random vector
  $\vb{X}^\tp=(X_1,X_2,\dots,X_p)$, and let $\Si$ have the
  eigenvalue--eigenvector pairs
  $(\la_1,\vb{e}_1),(\la_2,\vb{e}_2),\dots,(\la_p,\vb{e}_p)$ with
  $\la_1\ge\la_2\ge\dots\ge\la_p\ge 0$. Then the $i$-th principal component is
  given by
  \begin{multline}
    \label{eq:pc_definition}
    Y_i=\vb{e}_i^\tp\vb{X}\\
    =e_{i1}X_1+e_{i2}X_2+\dots+e_{ip}X_p,
  \end{multline}
  for $i=1,2,\dots,p$, and with these choices
  \begin{align}
    \Var(Y_i)     &=\vb{e}_i^\tp\Si\,\vb{e}_i=\la_i,
    \label{eq:pc_var}\\
    \Cov(Y_i,Y_k) &=\vb{e}_i^\tp\Si\,\vb{e}_k=0,\quad i\neq k.
    \label{eq:pc_cov}
  \end{align}
\end{theorem}

\begin{remark}
  If some $\la_i$ are equal, the choices of the corresponding coefficient
  vectors $\vb{e}_i$ --- and hence of $Y_i$ --- are not unique.
\end{remark}

\subsection{Uncorrelatedness of the components}
\label{subsec:pca_uncorrelated}

It remains to show that the components are uncorrelated, that is, that
$\vb{e}_i\perp\vb{e}_k \imp \vb{e}_i^\tp\vb{e}_k=0$ for $i\neq k$ gives
$\Cov(Y_i,Y_k)=0$.

The eigenvectors of $\Si$ are orthogonal if all the eigenvalues
$\la_1,\la_2,\dots,\la_p$ are distinct. If they are not all distinct, the
eigenvectors corresponding to common eigenvalues may be \emph{chosen} to be
orthogonal. Therefore, for any two eigenvectors $\vb{e}_i$ and $\vb{e}_k$ we
have $\vb{e}_i^\tp\vb{e}_k=0$ for $i\neq k$, and since
$\Si\vb{e}_k=\la_k\vb{e}_k$, premultiplying by $\vb{e}_i^\tp$ gives
\begin{equation}
  \label{eq:cov_zero}
  \Cov(Y_i,Y_k)=\vb{e}_i^\tp\Si\,\vb{e}_k
  =\vb{e}_i^\tp\la_k\vb{e}_k
  =\la_k\,\vb{e}_i^\tp\vb{e}_k=0,
\end{equation}
for all $i\neq k$.

\subsection{Total population variance}
\label{subsec:total_variance}

\begin{theorem}[Decomposition of the total variance]
  \label{thm:total_variance}
  Let $\vb{X}^\tp=(X_1,X_2,\dots,X_p)$ have covariance matrix $\Si$ with pairs
  $(\la_1,\vb{e}_1),\dots,(\la_p,\vb{e}_p)$, where
  $\la_1\ge\la_2\ge\dots\ge\la_p\ge0$, and let
  $Y_1=\vb{e}_1^\tp\vb{X},\dots,Y_p=\vb{e}_p^\tp\vb{X}$ be the principal
  components. Then
  \begin{multline}
    \label{eq:total_variance}
    \si_{11}+\si_{22}+\dots+\si_{pp}=\sum_{i=1}^{p}\Var(X_i)\\
    =\la_1+\la_2+\dots+\la_p=\sum_{i=1}^{p}\Var(Y_i).
  \end{multline}
\end{theorem}

\begin{proof}
  From the definition of the trace,
  $\si_{11}+\si_{22}+\dots+\si_{pp}=\tr(\Si)$. With $\mat{A}=\Si$ we may write
  the spectral decomposition
  \[
    \Si=\mat{P}\La\mat{P}^\tp,
  \]
  where $\La$ is the diagonal matrix of eigenvalues and
  $\mat{P}=(\vb{e}_1,\vb{e}_2,\dots,\vb{e}_p)$ is orthogonal, so that
  $\mat{P}^\tp\mat{P}=\mat{P}\mat{P}^\tp=\I$. Then, by the cyclic property of
  the trace,
  \[
    \tr(\Si)=\tr(\La\mat{P}^\tp\mat{P})=\tr(\La)
    =\la_1+\la_2+\dots+\la_p.
  \]
  Thus
  \[
    \sum_{i=1}^{p}\Var(X_i)=\tr(\Si)=\tr(\La)=\sum_{i=1}^{p}\Var(Y_i),
  \]
  because $\Var(Y_i)=\la_i$ by Theorem~\ref{thm:pca_main}.
\end{proof}

Consequently the \textbf{total population variance} is
\begin{equation}
  \label{eq:total_pop_variance}
  \si_{11}+\si_{22}+\dots+\si_{pp}=\la_1+\la_2+\dots+\la_p,
\end{equation}
and the proportion of it accounted for by the $k$-th component is
\begin{equation}
  \label{eq:proportion_variance}
  \qty{\begin{array}{c}
    \text{proportion of total}\\
    \text{variance due to the}\\
    \text{$k$-th component}
  \end{array}}
  =\frac{\la_k}{\la_1+\la_2+\dots+\la_p},
\end{equation}
for $k=1,2,\dots,p$.

\subsection{Interpretation of the loadings}
\label{subsec:loadings}

Consider each component of $\vb{e}_i=(e_{i1},e_{i2},\dots,e_{ip})$. The
magnitude of $e_{ik}$ measures the importance of the $k$-th variable to the
$i$-th principal component; in particular $e_{ik}$ is proportional to the
correlation coefficient between $Y_i$ and $X_k$,
\[
  e_{ik}\;\propto\;\Corr(Y_i,X_k).
\]

\begin{proposition}[Correlation of components and variables]
  \label{prop:corr_pc_var}
  If $Y_1=\vb{e}_1^\tp\vb{X},\dots,Y_p=\vb{e}_p^\tp\vb{X}$ are the principal
  components obtained from $\Si$, then
  \begin{equation}
    \label{eq:corr_pc_var}
    \rho_{Y_i,X_k}=\frac{e_{ik}\sqrt{\la_i}}{\sqrt{\si_{kk}}},
    \qquad i,k=1,\dots,p.
  \end{equation}
\end{proposition}

\subsection{Principal components from standardized variables}
\label{subsec:standardized}

The standardized variables are
\begin{equation}
  \label{eq:standardized_vars}
  Z_1=\frac{X_1-\mu_1}{\sqrt{\si_{11}}},\;\dots,\;
  Z_p=\frac{X_p-\mu_p}{\sqrt{\si_{pp}}}.
\end{equation}
In matrix notation, where $\mat{V}^{1/2}$ is the diagonal standard deviation
matrix,
\begin{equation}
  \label{eq:standardized_matrix}
  \vb{Z}=\qty{\mat{V}^{1/2}}\inv(\vb{X}-\mub),
  \qquad \E(\vb{Z})=\vb{0},
\end{equation}
and
\begin{equation}
  \label{eq:cov_Z_is_rho}
  \Cov(\vb{Z})=\qty{\mat{V}^{1/2}}\inv\Si\qty{\mat{V}^{1/2}}\inv=\rhob .
\end{equation}
The principal components may therefore be obtained from the eigenvectors of the
\emph{correlation matrix} $\rhob$ of $\vb{X}$. All the previous results apply,
with some simplification given that the variance of each $Z_i$ is unity.

\begin{theorem}[Components of the standardized variables]
  \label{thm:pca_standardized}
  The $i$-th principal component of the standardized variables
  $\vb{Z}^\tp=(Z_1,Z_2,\dots,Z_p)$ with $\Cov(\vb{Z})=\rhob$ is given by
  \begin{equation}
    \label{eq:pc_standardized}
    Y_i=\vb{e}_i^\tp\vb{Z}
       =\vb{e}_i^\tp\qty{\mat{V}^{1/2}}\inv(\vb{X}-\mub),
  \end{equation}
  where $(\la_1,\vb{e}_1),\dots,(\la_p,\vb{e}_p)$ are the
  eigenvalue--eigenvector pairs of $\rhob$ with
  $\la_1\ge\la_2\ge\dots\ge\la_p\ge0$. Moreover
  \begin{align}
    \sum_{i=1}^{p}\Var(Y_i)&=\sum_{i=1}^{p}\Var(Z_i)=p,
    \label{eq:standardized_total_variance}\\
    \rho_{Y_i,Z_k}&=e_{ik}\sqrt{\la_i}.
    \label{eq:corr_standardized}
  \end{align}
\end{theorem}

We see from \eqref{eq:standardized_total_variance} that the total standardized
population variance is simply $p$, the sum of the diagonal elements of the
matrix $\rhob$. Working with $\vb{Z}$ in place of $\vb{X}$, we find the
proportion of total variance explained by the $k$-th principal component of
$\vb{Z}$ to be
\begin{equation}
  \label{eq:proportion_standardized}
  \qty{\begin{array}{c}
    \text{proportion of standardized}\\
    \text{variance due to the}\\
    \text{$k$-th component}
  \end{array}}
  =\frac{\la_k}{p},
\end{equation}
for $k=1,2,\dots,p$, where the $\la_k$ are now the eigenvalues of $\rhob$.

\selectlanguage{spanish}
\subsection{An\'alisis de componentes principales}
\label{subsec:acp_espanol}

Sea $\Si$ la matriz de covarianzas del conjunto de datos. La varianza total se
da por la traza
\begin{equation}
  \label{eq:varianza_total_es}
  \text{varianza total}=\tr\Si=\sum_{i=1}^{p}\Si_{ii},
\end{equation}
donde $\Si_{ii}$ son las varianzas individuales de cada variable.

Usando la descomposici\'on en valores y vectores propios (SVD) de $\Si$, existen
matrices $\mat{Q}$ y $\La$ tales que
\begin{equation}
  \label{eq:svd_es}
  \Si=\mat{Q}\La\mat{Q}^\tp,
\end{equation}
con $\La$ diagonal con valores propios $\la_1,\dots,\la_p$. La varianza total
puede escribirse como
\begin{multline}
  \label{eq:traza_es}
  \text{varianza total}=\tr\Si=\tr\qty{\mat{Q}\La\mat{Q}^\tp}\\
  =\tr\La=\la_1+\la_2+\dots+\la_p.
\end{multline}

Las columnas de $\mat{Q}$ son las direcciones de las componentes principales;
los valores de $\La$ indican cu\'anta varianza explica cada componente. Estas
componentes definen un nuevo sistema de coordenadas: cada punto se representa
por sus puntuaciones de componentes, relacionadas con las variables originales
por
\begin{equation}
  \label{eq:reconstruccion_es}
  \vb{X}=\mat{Q}(\mathrm{PC})+\bar{\vb{x}},
\end{equation}
donde $\vb{X}$ es el vector de coordenadas originales, $\mathrm{PC}$ el de
puntuaciones y $\bar{\vb{x}}$ el vector de medias.

\subsection{Resumen del proceso de PCA}
\label{subsec:resumen_pca}

Se resume seg\'un estos pasos:
\begin{enumerate}
  \item \textbf{Estandarizar los datos.}
  \item \textbf{Calcular la matriz de covarianzas.}
    \begin{itemize}
      \item La matriz calcula c\'omo var\'ian conjuntamente las variables.
      \item Su tama\~no es $p\times p$, donde $p$ es el n\'umero de variables.
    \end{itemize}
  \item \textbf{Obtener la descomposici\'on en valores y vectores propios
        (SVD).}
    \begin{itemize}
      \item Los \emph{eigenvalues} miden la varianza explicada por cada
            componente.
      \item Los \emph{eigenvectors} son las direcciones (combinaciones
            lineales) de las componentes.
    \end{itemize}
  \item \textbf{Ordenar las componentes.}
    \begin{itemize}
      \item Ordenar valores propios de mayor a menor.
      \item El vector correspondiente al mayor valor propio es la primera
            componente.
    \end{itemize}
  \item \textbf{Decidir cu\'antas componentes conservar.}
    \begin{itemize}
      \item Criterio de Kaiser: conservar valores propios $>1$.
      \item Gr\'afico \emph{scree}: buscar el ``codo''.
      \item Varianza acumulada.
    \end{itemize}
  \item \textbf{Obtener puntuaciones:} proyectar los datos centrados sobre las
        componentes seleccionadas,
        \begin{equation}
          \label{eq:puntuaciones_es}
          \mathrm{PC}=\mat{Q}^\tp\qty{\vb{X}-\bar{\vb{x}}}.
        \end{equation}
\end{enumerate}
\selectlanguage{english}
